\documentclass[../../../include/open-logic-section]{subfiles}

\begin{document}

\olfileid[id]{sth}{z}{pairs}
\olsection{Pasangan}

Aksioma berikutnya yang akan kita pertimbangkan adalah sebagai berikut:

\begin{axiom}[Pasangan]
Untuk sebarang himpunan $a,b$, himpunan $\{a,b\}$ ada.
\[
	\forall a \forall b \exists P \forall x (x \in P \liff (x = a \lor x = b))
\]
\end{axiom}

Berikut cara membenarkan aksioma ini dengan menggunakan konsepsi iteratif.
Andaikan $a$ tersedia pada tahap~$S$, dan $b$ tersedia pada tahap~$T$.
Misalkan $M$ adalah tahap yang lebih akhir di antara $S$ dan~$T$. Karena $a$
dan $b$ sama-sama tersedia pada tahap~$M$, himpunan $\{a,b\}$ merupakan
kumpulan yang mungkin dan tersedia pada setiap tahap sesudah~$M$.

Namun, tunggu!{} Mengapa kita mengandaikan bahwa memang \emph{ada} tahap sesudah
$M$? Jika tidak ada, pembenaran kita gagal. Jadi, untuk membenarkan Pasangan,
kita harus menambahkan satu prinsip lagi pada kisah yang disampaikan dalam
\olref[sth][z][story]{sec}, yakni:
\begin{enumerate}
	\item[] \stagessucc. Tidak ada tahap terakhir.
\end{enumerate}
Apakah prinsip ini dapat dibenarkan? Tidak ada bagian kisah Shoenfield yang
menyatakan \emph{secara eksplisit} bahwa tidak ada tahap terakhir. Meskipun
demikian, sekalipun prinsip itu---secara tegas---merupakan tambahan pada kisah
kita, prinsip tersebut sangat selaras dengan gagasan dasar bahwa himpunan
dibentuk secara bertahap. Untuk selanjutnya, kita akan menerimanya. Dengan
demikian, kita juga menerima Aksioma Pasangan.

Berbekal aksioma baru ini, kita dapat membuktikan keberadaan jauh lebih banyak
himpunan. Sebagai contoh:

\begin{prop}\ollabel{prop:pairsconsequences}
Untuk sebarang himpunan $a$ dan~$b$, himpunan-himpunan berikut ada:
	\begin{enumerate}
		\item\ollabel{singleton} $\{a\}$
		\item\ollabel{binunion} $a \cup b$
		\item\ollabel{tuples} $\tuple{a,b}$
	\end{enumerate}
\end{prop}

\begin{proof}
\olref{singleton}. Berdasarkan Pasangan, $\{a,a\}$ ada; menurut
Ekstensionalitas, himpunan ini sama dengan $\{a\}$.

\olref{binunion}. Berdasarkan Pasangan, $\{a,b\}$ ada. Maka
$a\cup b=\bigcup\{a,b\}$ ada berdasarkan Gabungan.

\olref{tuples}. Berdasarkan \olref{singleton}, $\{a\}$ ada. Berdasarkan
Pasangan, $\{a,b\}$ ada. Maka
$\{\{a\},\{a,b\}\}=\tuple{a,b}$ ada, sekali lagi berdasarkan Pasangan.
\end{proof}

\begin{prob}
Tunjukkan bahwa, untuk sebarang himpunan $a,b,c$, himpunan $\{a,b,c\}$ ada.
\end{prob}

\begin{prob}
Tunjukkan bahwa, untuk sebarang himpunan $a_1,\ldots,a_n$, himpunan
$\{a_1,\ldots,a_n\}$ ada.
\end{prob}

%\begin{proof} By two applications of Pairs, $\{\{a_1, a_2\}, \{a_1,
%   a_3\}\}$ exists. By Union and Extensionality, $\bigcup \{\{a_1,
%   a_2\}, \{a_1, a_3\}\} = \{a_1, a_2, a_3\}$ exists. Repeat this
%   trick as often as necessary. \end{proof}

\end{document}
